\thispagestyle{empty}
\vspace*{12mm}
{\color{accent}\sffamily\large SOURCE-PINNED PROGRAMMER'S EDITION\par}
\vspace{12mm}
{\sffamily\bfseries\fontsize{31}{35}\selectfont DSLtk and MetaML\par}
\vspace{8mm}
{\sffamily\fontsize{21}{26}\selectfont A Full Guide for the Working Programmer\par}
\vspace{15mm}
{\Large C++20 and D\par}
\vspace{6mm}
A practical guide to constructing languages, embedding MetaML, and validating the boundary between scripts and native applications.
\vfill
24 chapters\quad /\quad 120 paired workshops\\[4mm]
Two complete case studies\quad /\quad 244 native example programs
\vspace{10mm}
\clearpage
\section*{About this edition}
This guide is based on the four source files supplied with this request: DSLtk(3).hpp, DSLtk(1).d, MetaML(2).hpp and MetaML.d. Canonical copies live in vendor/ as DSLtk.hpp, DSLtk.d, MetaML.hpp and MetaML.d. SOURCE-MAP.md records SHA-256 hashes, original names and important source locations. The guide describes this snapshot; it is not an independently versioned upstream specification.

The C++ headers require C++20. D uses template mixins, modules and Phobos facilities rather than literal header inclusion. The examples deliberately preserve native idioms where the host languages differ. MetaML syntax inside source strings is a third language and should not be confused with either host.

Every workshop has one C++ program and one D program. The four standalone case-study programs bring the total to 244. The C++ examples are compiled and executed against the supplied snapshot. No D compiler is installed in the authoring environment, so D code is source-reviewed and supplied with a runnable test target, not described as compiler-validated.

The external MetaML standard-library tree was not supplied. Examples therefore use syntax, constructor declarations and core primitives available in a fresh session. Standard-library loading is explained as a separate installation step, with missing-library behavior tested explicitly.

\clearpage
\section*{Who this book is for}
You should be comfortable writing functions, structs or classes, using collections and compiling a native program. The C++ chapters assume basic familiarity with templates and lambda captures. The D chapters assume familiarity with imports, delegates and template instantiation. Advanced host metaprogramming is introduced only where it changes how the library is used.

Read the first four chapters even if your immediate goal is embedding MetaML. They establish ownership, value semantics and error channels that reappear at every native boundary. Then follow one of two routes: chapters 5 through 12 for constructing a language with DSLtk, or chapters 13 through 22 for using MetaML as a ready-made language.

Each workshop occupies a two-page spread: the C++ implementation and execution reasoning, then the D implementation, portability notes and failure-oriented exercise. These spreads are intentionally small enough to run independently. A chapter opening provides the larger design context, and its final page gives implementation notes and a review path.

The case studies at the end combine the facilities into a formula optimizer and a reloadable policy core. They include complete native programs, negative cases and acceptance checks. They are synchronous tools with deliberately defined scope; neither claims to be a production server or a general-purpose algebra system.

\clearpage
\section*{How to use the examples}
Example IDs encode chapter and workshop number. ex05\_03 is the third workshop in chapter 5. A complete C++ translation unit lives in examples/cpp/ex05\_03.cpp and its D counterpart in examples/d/ex05\_03.d. The book lists the function body to avoid repeating imports on every page. The adjacent .inc files are the exact body listings used by LaTeX.

\begin{lstlisting}[language=C++]
#include "MetaML.hpp"
#include <cassert>
int main() {
  // Insert a C++ workshop body here.
}
\end{lstlisting}
\begin{lstlisting}[language=D]
import DSLtk;
import MetaML;
void main() {
  // Insert a D workshop body here.
}
\end{lstlisting}
Compile one example with the vendor include or import directory. For D, include both module sources when linking. examples/all\_cpp.cpp combines the 120 workshop bodies into a fast aggregate test; each individual file remains independently compilable. Optional Python tools regenerate example files from the editorial content data.

Assertions state the intended outcome. A negative example succeeds as a test only when the expected library failure occurs. Run examples with assertions enabled. Do not compile with NDEBUG or a D release configuration that removes assertions when validating the book.

\clearpage
\section*{DSLtk: a facility map}
Composition facilities include Pipeline, Operators, PatternMatch, AST, Rewrite, ExprTemplates, CustomLiterals and EnumReflection. These help define a native authoring surface. Their feature tags can be composed, while public factories remain callable without a facade.

Control and construction facilities include Maybe, Result, Lazy, memoize, parser combinators, task chains and semantic production handlers. Choose the smallest abstraction that matches your failure and lifetime requirements. A function pipeline and a task pipeline are not automatically asynchronous execution systems.

Language facilities include PEGDefinition, SExpr, EGraph, serialization protocols, abstract I/O, visitors, static byte regexes, BinaryStream, Preprocessor and CFG. Graphs, terms, unification, term rewriting and shared graph reduction support more explicit transformation systems.

Source infrastructure includes Span, Trivia, BasicToken, SourceMap and Diagnostic. The generic REPL Session is policy-based and distinct from metaml::Session. It expects language-specific parse and evaluation operations. MetaML provides ReplLanguage as an adapter.

The workshops emphasize executable entry points; implementation notes explain the remaining extension surfaces and their limits. When building a new DSL, keep a small representation and an explicit semantic contract rather than enabling every available feature at once.

\clearpage
\section*{MetaML: an embedding map}
The native embedding begins with Session. Its constructor installs core primitives and the Option and Result declarations. eval returns a Result containing the final Value or a vector of diagnostics. type\_of parses and checks an expression without evaluating its value. load\_file and run\_file provide file-oriented execution.

The language implementation includes lexing, parsing, an expression and item AST, type inference, module checking and an interpreter. Values include numbers, strings, lists, tuples, records, constructors, functions, modules and foreign objects. Host conversion adapters cover a defined set of native types.

The type checker supports Hindley-Milner-style inference, numeric constraints and row-polymorphic records. Modules and signatures organize public interfaces. Functional, tacit and concatenative forms share an evaluator, but terse syntax is best introduced after an equivalent explicit lambda is understood.

register\_primitive exposes trusted host functions; install\_primitive and the allowlist support more controlled extension installation. set\_output redirects session text. on\_close registers resource teardown, and close stops all resources before waiting. spawn\_interp has a separate restricted background-use design that must not be mistaken for blanket thread safety.

The external metalib tree supplies additional modules. Its discovery is explained, but no unprovided module is treated as bundled functionality. Portable examples use the same MetaML source in both native hosts.

\clearpage
\section*{Four error boundaries}
First, native build errors indicate a host compiler, language-standard, import-path or template-instantiation problem. Resolve these before interpreting script behavior. The build instructions provide a minimal command for each host and preserve the exact vendor snapshot.

Second, syntax failures mean the language text cannot produce a valid AST. Parser results identify incomplete or malformed input. A REPL may ask for more input for an incomplete form, while a batch runner must eventually report a final syntax error.

Third, checking failures include undefined names, incompatible types and invalid module interfaces. They should reach the user as structured diagnostics. Preserve codes and spans separately from formatted caret text, especially in editor or service integrations.

Fourth, execution failures can occur after a program has passed the checker. Native callbacks may fail, a script may hit a depth bound, or an explicit language error may be raised. Earlier committed units and external effects are not automatically rolled back.

There is also a domain outcome inside a successful language result: a policy can return false or a language-level Err constructor. This is not the same as a host Result error. Case study II keeps evaluation failure distinct from allow and deny. Record which boundary failed before choosing a user-facing response.

\clearpage
\section*{Syntax conventions}
C++ examples use dsl:: and metaml:: qualification. Compile-time string arguments use angle brackets, such as leaf<"number">. D uses template arguments such as leaf!"number" and dot qualification for enum members. Predicate negation uses ! in C++ and \textasciitilde{} or negated in D. Parser repetition uses unary star in C++ and many in D.

MetaML source appears inside native strings or in .ml fixtures. A value declaration uses val name := expression. Anonymous functions use fn parameters -> body. A local expression uses let name = expression in body. Modules use module Name := struct ... end. Record fields use equals, while the top-level declaration marker is colon-equals.

Escaping belongs to the host string layer: a newline in MetaML source may be written as backslash-n in a C++ or D string. A quote within MetaML string syntax must survive host string decoding. Prefer a separate .ml file for long scripts, and typed native primitives for runtime data.

Identifiers such as ?x inside rewrite S-expressions are rule variables. They are unrelated to MetaML type variables, which use an apostrophe in type syntax. A concrete atom x in an e-graph is data; an unbound ?x pattern cannot be inserted as a concrete expression.

Byte positions are zero-based offsets; diagnostic line and column metadata is one-based. Half-open spans include the first byte and exclude the end byte.

\clearpage
\section*{A working development loop}
Begin with one accepted input and one rejected input. Write down the output type, ownership and failure policy before adding syntax. Then create a minimal example that exercises exactly that boundary. A small test with an observable result is more useful than a large demonstration that silently ignores errors.

For a new DSL, choose a representation and validate its invariants before writing optimizations. Keep recognition separate from semantic construction when actions might otherwise escape a failed parse. For rewrites, state the integer model, effects model and traversal strategy, and record the stop reason.

For embedding MetaML, bring up a fresh session without the external prelude first. Add typed host primitives one at a time. Replace the final allowlist after registration, configure output and resource shutdown, and prepare source in a candidate environment before publication.

For cross-dialect work, run identical .ml fixtures in both hosts and compare values, types, diagnostic codes and spans. Native adapter tests should remain separate because they exercise host conversion and lifetime behavior that identical scripts cannot reveal.

The book test runner checks the C++ workshop programs and complete case studies. Its report is included in VALIDATION.md. A missing D toolchain is explicitly reported. A reproducible build should distinguish completed checks from checks that remain to be run on another machine.

\clearpage
\section*{Building and navigating}
The normal build requires a TeX distribution with LuaLaTeX and the packages used by book.cls. Run texlua build.lua from this directory, or use make. latexmk -lualatex main.tex is the alternative route configured by .latexmkrc. All routes create build/main.pdf and preserve source files.

The supplied PDF has 300 physical pages. The layout uses explicit workshop spreads, not blank padding. There are twelve opening pages followed by twenty-four chapters of twelve pages each. Editing prose or fonts can change pagination; the check target validates the page count for the shipped edition rather than hiding overflow.

The table of contents on the next two pages links directly to chapter starts. Chapter openings give the design context. Each pair of workshop pages shares an example ID, and every example has a source anchor into vendor/. The final page of each chapter collects extension notes and review questions.

BUILD.md documents dependencies, all targets, cleanup and per-example commands. SOURCE-MAP.md identifies the input snapshots. VALIDATION.md records actual test evidence. The ZIP includes the compiled PDF as a convenience, all LaTeX source, every example and complete case-study programs.

No shell escape is required to build the book. Code listings use the listings package, so there is no Pygments dependency. Python is optional for ordinary TeX builds and is used by development and validation commands.

\clearpage
\section*{Contents}
Chapter starts are linked. Page numbers refer to the physical pagination of this edition.

\noindent\hyperref[ch:1]{\textbf{1.} The two libraries and the two host languages}\hfill 13\par\medskip
\noindent\hyperref[ch:2]{\textbf{2.} Pipelines, predicates and dispatch}\hfill 25\par\medskip
\noindent\hyperref[ch:3]{\textbf{3.} AST values and representation contracts}\hfill 37\par\medskip
\noindent\hyperref[ch:4]{\textbf{4.} Absence, errors and deferred work}\hfill 49\par\medskip
\noindent\hyperref[ch:5]{\textbf{5.} Parser combinators and committed failure}\hfill 61\par\medskip
\noindent\hyperref[ch:6]{\textbf{6.} PEG-style token grammars}\hfill 73\par\medskip
\noindent\hyperref[ch:7]{\textbf{7.} S-expressions and rewrite rules}\hfill 85\par\medskip
\noindent\hyperref[ch:8]{\textbf{8.} Graphs, terms and bounded reduction}\hfill 97\par\medskip
\noindent\hyperref[ch:9]{\textbf{9.} Equality saturation in practice}\hfill 109\par\medskip
\noindent\hyperref[ch:10]{\textbf{10.} Byte recognition, binary streams and CFGs}\hfill 121\par\medskip
\noindent\hyperref[ch:11]{\textbf{11.} Preprocessing and diagnostics}\hfill 133\par\medskip
\noindent\hyperref[ch:12]{\textbf{12.} Visitors, literals and reflection}\hfill 145\par\medskip
\clearpage
\section*{Contents / continued}
Chapter starts are linked. Page numbers refer to the physical pagination of this edition.

\noindent\hyperref[ch:13]{\textbf{13.} MetaML evaluation and session boundaries}\hfill 157\par\medskip
\noindent\hyperref[ch:14]{\textbf{14.} Numbers, strings and collections}\hfill 169\par\medskip
\noindent\hyperref[ch:15]{\textbf{15.} Functions, application and lexical scope}\hfill 181\par\medskip
\noindent\hyperref[ch:16]{\textbf{16.} Type inference and records}\hfill 193\par\medskip
\noindent\hyperref[ch:17]{\textbf{17.} Algebraic data types and pattern matching}\hfill 205\par\medskip
\noindent\hyperref[ch:18]{\textbf{18.} Modules and explicit interfaces}\hfill 217\par\medskip
\noindent\hyperref[ch:19]{\textbf{19.} Typed native primitives}\hfill 229\par\medskip
\noindent\hyperref[ch:20]{\textbf{20.} Diagnostics and inspection tools}\hfill 241\par\medskip
\noindent\hyperref[ch:21]{\textbf{21.} Capabilities, callbacks and shutdown}\hfill 253\par\medskip
\noindent\hyperref[ch:22]{\textbf{22.} Deployment and cross-dialect verification}\hfill 265\par\medskip
\noindent\hyperref[ch:23]{\textbf{23.} Case study I: a formula optimization tool}\hfill 277\par\medskip
\noindent\hyperref[ch:24]{\textbf{24.} Case study II: a reloadable decision service}\hfill 289\par\medskip
